\documentclass[10pt]{article}
\usepackage{times}
\usepackage[margin=1in]{geometry}
\begin{document}
\section*{Transmission of signals}
The message may be a function $f(x, y, t)$ of two space coordinates and time,
the light intensity at point $(x, y)$ and time $t$ on a pickup tube plate. Two or
more functions of time, say $f(t)$, $g(t)$, $h(t)$ --- this is the case in
three dimensional sound transmission or if the system is intended to service
several individual channels in multiplex. The logarithmic measure is more
convenient: $\log_2 M = \log_{10} M / \log_{10} 2 = 3.32 \log_{10} M$, and a
device with two stable positions, such as a relay or a flip-flop circuit, can
store one bit of information. $N$ such devices can store $N$ bits, since the
total number of possible states is $2^N$ and $\log_2 2^N = N$.
If the source has $n$ symbols with probabilities $(p_1, p_2, \ldots, p_n)$
then the entropy $H = -\sum_{i} p_i \log p_i$ measures, in bits per symbol,
the information produced. Several functions $(u, v, w)$ of several variables
$(x, y, z, t)$ arise in colour television, and the channel capacity
$C = \lim_{T\to\infty} (\log N(T))/T$ bounds the rate of transmission.
\end{document}
